\documentclass[11pt]{article}
\usepackage[T1]{fontenc}
\usepackage{amsmath,amssymb,amsthm}
\usepackage{geometry}
\geometry{margin=1in}
\newtheorem{theorem}{Theorem}
\newtheorem{corollary}{Corollary}
\newtheorem{proposition}{Proposition}
\newcommand{\pw}{\operatorname{pw}}
\newcommand{\lbw}{\operatorname{lbw}}
\title{Laboratory V71: Support-Hypergraph Linear Branch-Width and Primal Pathwidth}
\author{NC0\_k-Avoid Laboratory}
\date{31 July 2026}
\begin{document}
\maketitle

\paragraph{Status.}
This is an internally verified theorem module. It is not peer reviewed, novelty confirmed, or a claim about resolving P versus NP.

Let $H=(X,E)$ be a finite hypergraph of rank at most $r$.  For $S\subseteq E$, define
\[
  V(S)=\bigcup_{e\in S}e,
  \qquad
  \lambda_H(S)=|V(S)\cap V(E\setminus S)|.
\]
For an ordering $\pi=(e_1,\ldots,e_m)$, put $S_i=\{e_1,\ldots,e_i\}$ and
\[
  q(\pi)=\max_i\lambda_H(S_i),
  \qquad
  q^*(H)=\min_\pi q(\pi).
\]
Let $P(H)$ be the primal graph on $X$, with two variables adjacent when they occur together in a hyperedge.

\begin{theorem}[Exact width vocabulary]
The V70 support frontier at $S$ has cardinality $\lambda_H(S)$.  Consequently $q^*(H)$ is exactly the linear branch-width of the hyperedge ground set under the vertex-boundary connectivity function $\lambda_H$.  For rank-two hypergraphs this specializes to the classical graph linear-width edge layout.
\end{theorem}

\begin{proof}
The V70 frontier is
\[
F(S)=\left(\bigcup_{e\in S}e\right)\cap
     \left(\bigcup_{e\notin S}e\right),
\]
which is precisely the boundary whose size is $\lambda_H(S)$.  A linear branch decomposition is equivalently a linear ordering of the ground-set elements, and each cut of the caterpillar induces one prefix--suffix partition.  Thus the layout width is $q(\pi)$ and its optimum is $q^*(H)$.
\end{proof}

\begin{theorem}[Pathwidth sandwich]
For every rank-at-most-$r$ hypergraph,
\[
  q^*(H)\le \pw(P(H))+1
  \qquad\text{and}\qquad
  \pw(P(H))\le q^*(H)+r-1.
\]
In particular, for ternary supports,
\[
  q^*(H)\le \pw(P(H))+1,
  \qquad
  \pw(P(H))\le q^*(H)+2.
\]
Both transformations are constructive.
\end{theorem}

\begin{proof}
First fix an edge order $\pi=(e_1,\ldots,e_m)$ and write $S_i=\{e_1,\ldots,e_i\}$.  Define
\[
  B_i=F(S_{i-1})\cup e_i.
\]
Every primal edge lies in some hyperedge and hence in the corresponding bag.  If a variable $x$ first occurs at position $a$ and last occurs at position $b$, then $x\in e_a\subseteq B_a$, $x\in e_b\subseteq B_b$, and for every $a<i<b$ it occurs on both sides of the cut, so $x\in F(S_{i-1})\subseteq B_i$.  Therefore the bags containing $x$ form an interval.  The sequence is a path decomposition and $|B_i|\le q(\pi)+r$.  Hence $\pw(P(H))\le q(\pi)+r-1$ and minimization proves the second inequality.

Conversely let $(B_1,\ldots,B_t)$ be a path decomposition of width $p$.  Every hyperedge is a clique of $P(H)$.  The bag-index intervals of its vertices pairwise intersect, and intervals satisfy the Helly property, so some bag contains the whole hyperedge.  Let $\rho(e)$ be the rightmost such bag and sort hyperedges by nondecreasing $\rho$, breaking ties arbitrarily.  Consider any cut of this order and any boundary variable $x$.  A processed support containing $x$ has a complete-support bag no later than the current $\rho$ value, while an unprocessed support containing $x$ has one no earlier.  The bags containing $x$ form an interval, so $x$ belongs to the current bag.  For a cut inside a tie block, both relevant supports are contained in the tied bag.  Thus every boundary has size at most $p+1$, proving $q^*(H)\le p+1$.
\end{proof}

Let $A(b)$ denote the number of nonempty affine subspaces of $\mathbb F_2^b$.

\begin{corollary}[Constructive projected-DAG bound]
If a path decomposition of $P(H)$ of width $p$ is supplied, then a gate order can be constructed in polynomial time such that
\[
  G_{\mathrm{proj}}\le m A(p+1).
\]
\end{corollary}

\begin{proof}
The rightmost-bag construction gives an order $\pi$ with $q(\pi)\le p+1$.  V70 proves $G_{\mathrm{proj}}(\pi)\le mA(q(\pi))$, and $A$ is monotone.
\end{proof}

\begin{proposition}[Tree-decomposition affine feasibility DP]
Given a nice tree decomposition of $P(H)$ of width $k$, affine-cell consistency can be decided by a bottom-up dynamic program storing at most $A(k+1)$ distinct nonempty affine residuals per bag.  A direct join examines at most $A(k+1)^2$ state pairs, with polynomial-time linear algebra in $k$.
\end{proposition}

\begin{proof}
Assign every gate to a bag containing its complete support.  Such a bag exists because the support is a clique in the primal graph and vertex subtrees in a tree decomposition obey the Helly property.  At an assigned gate, intersect each incoming residual with either affine cell.  At a forget node, existentially project the forgotten coordinate; at an introduce node, add an unconstrained coordinate; at a join, intersect one residual from each child and discard inconsistent results.  Deduplicate canonical affine subspaces after every operation.  Since a bag contains at most $k+1$ variables, no table contains more than all nonempty affine subspaces of $\mathbb F_2^{k+1}$, namely $A(k+1)$.  The direct join bound follows.
\end{proof}

\paragraph{Boundary of the result.}
The tree-decomposition DP is tree-shaped and is not the linear projected residual DAG of V68--V70.  Treewidth alone cannot bound the linear layout parameter: trees have treewidth one but unbounded pathwidth.  No unrestricted polynomial good-order theorem, all-orders lower bound, standard proof-system simulation, unrestricted range-avoidance algorithm, or P-versus-NP consequence follows.

\end{document}
